\Acp{VLM} in \ac{RS} enable natural language interactions with images across multiple tasks, including \ac{VQA}, change captioning, and \genDet{} \cite{kuckreja2024geochat, irvin2025teochat, soni2025earthdial}. To facilitate the research and development of \acp{VLM}, several image-text datasets have been proposed in RS.
Early \ac{RS} image–text datasets, such as UCM-Captions, Sydney-Captions \cite{qu2016deep}, RSICD \cite{lu2017exploring}, RSITMD \cite{yuan2022exploring}, and NWPU-Captions \cite{cheng2022nwpu} contain high-resolution aerial imagery with human-written single-sentence descriptions with limited intra-class diversity. 
To overcome this bottleneck, recently proposed datasets shift toward large-scale, automatically generated texts: RS5M \cite{zhang2024rs5m} introduces the first million-scale corpus with 5 million image–caption pairs; GAIA \cite{gaia_zavras} provides \num{205150} pairs, with multi-sentence captions emphasizing environmental dynamics; and Landsat30-AU \cite{ma2025landsat30} contains \num{196262} image-caption pairs and \num{17725} \ac{MC} \acp{VQA}.
For \ac{VQA} specifically, RSVQA \cite{lobry2020rsvqa} comprises \num{77232} and \num{1066316} question-answer triplets for low- and high-resolution imagery, while RSVQAxBEN \cite{lobryRSVQAMeetsBigearthnet2021} scales to approximately \num{15} million pairs with \ac{S2} BigEarthNet~\cite{sumbul2019bigearthnet} \ac{RGB} composites for \ac{CLC} class presence.
Recent studies also extend the spectral range of image-text datasets beyond \ac{RGB}. 
Yuan et. al~\cite{ChatEarthNet} generate captions for nine-band \ac{S2} images, yet descriptions rely on broad \ac{LULC} classes from ESA WorldCover~\cite{zanaga2022esaworldcover} (\eg water, tree, snow). Marimo et al.~\cite{marimo2025beyond} scale to \num{975000} co-registered \ac{S1}–\ac{S2} images from SSL4EO~\cite{wang2023ssl4eo} but generate captions using a \ac{VLM} that receives \ac{RGB} bands only, preventing texts from relying on information beyond the visible spectrum. 